\clearpage
\fancyhead[L]{\small BEER AND MÄDLER.}
\hypertarget{native-657-lower}{}
\begin{center}\small Göttingische gelehrte Anzeigen. No. 137. pp. 1361..1365. 30 August 1841.\\[3pt]\rule{.72\linewidth}{.3pt}\end{center}

\emph{Contributions to Knowledge of the Physical Nature of the Celestial Bodies in the Solar System by} \textsc{Wilhelm Beer}, \emph{Royal Prussian Privy Councillor of Commerce, and Dr.} J. H. \textsc{Mädler}, \emph{Imperial Russian Court Councillor.} Weimar 1841. From B. F. Voigt. 152 pages in quarto, with 7 copperplates.

\emph{Fragments sur les corps célestes du système solaire par} \textsc{Guillaume Beer} \emph{et} \textsc{Jean Henri Mädler.} Paris 1840. Chez Bachelier. 216 pages in quarto, with 7 plates.

From the same hands to which we owe the classic selenography that appeared four years ago, we now receive a collection of short essays in German and French, parts of some of which had already been made known. They deal chiefly with the physical nature of the Moon and the planets. Bearing the marks of careful observation, rigorous examination and sound judgment, they extend our knowledge in several valuable ways. The subjects of the individual essays are as follows.
I. \emph{On the far hemisphere of the Moon.} The half of the Moon's surface turned away from us falls into two parts, one that becomes visible to us under favourable conditions of libration, and another that can never be seen from the Earth. This second part is described here as bounded by an ellipse with double lateral flattening, and its area is given as 0.4243 of the entire lunar surface. The reviewer finds that its shape is a \emph{quadrilateral} and its area 0.4118. Of course, we can never obtain any positive knowledge of the physical nature of this part. The authors rightly observe, however, that there is no reason at all to suspect any substantial difference in that respect between the two halves of the Moon, especially since the part of the far half that becomes visible through libration shows formations very similar to those on the near half. The only difference between the two halves that we know with certainty is the difference in their illumination. The authors set out these conditions in great detail, going so far as to describe the advantages for observational astronomy that the far hemisphere of the Moon would have over other sites if the astronomer were to imagine himself there.
II. \emph{On rilles in the lunar surface.} After \textsc{Schröter}, who 50 \SourceLeaf{658}years ago first discovered the two rilles near Hyginus and Ariadäus, a number of others were observed by \textsc{Pastorf}, \textsc{Gruithuysen} and especially \textsc{Lohrmann}. Their total number, however, is insignificant compared with the great many similar formations the authors found while preparing their map of the Moon. The list given here contains 92 of them.
III. \emph{On lunar eclipses.} Formulas for calculating them, particularly with regard to the enlargement that must be applied to the radius of the shadow cone to bring the observations into agreement with the calculation. Applying the formulas to three eclipses, on 26 December 1833, 10 June 1835 and 13 October 1837, gives very unequal results, the first $\frac{1}{65.4}$, the second $\frac{1}{28.3}$, the third $\frac{1}{54.0}$. The phenomenon is usually attributed to the Earth's atmosphere, which allows only a small proportion of the Sun's rays to pass through and thus increases the diameter of the opaque body of the Earth and hence also of the shadow cone. But this explanation is in any event very unsatisfactory when one looks more closely at the quantitative aspects of the phenomenon. Strong absorption of light can be conceded only for the lower, denser layers of air, while the light passing through the higher layers is attenuated less and less and at the same time, by refraction, actually reduces the diameter of the shadow cone. To explain the phenomenon observed in the lunar eclipse of 10 June 1835 by this cause, one would have to assume an opaque atmosphere 23 miles high. Thus, while allowing that this cause may contribute, we must admit that observers of a lunar eclipse believe they already see a point on the lunar disk enter the full shadow while it is in fact still receiving light from a small part of the solar disk. In the case in question, the calculation shows that, if attenuation by the Earth's atmosphere is disregarded, $\frac{1}{43}$ of the solar disk was still uncovered as seen from points on the apparent boundary of the umbra. Under the same restriction, the figure is $\frac{1}{178}$ for the eclipse of 1833 and $\frac{1}{126}$ for that of 1837.
Essays IV and V contain valuable monographs on interesting lunar landscapes, the region of the crater \textsc{Schröter} and the vicinity of the Moon's north pole. Essay V concerns Saturn and contains a very worthwhile attempt to determine the orbits of its two inner satellites from a new reduction of \textsc{Herschel}'s observations of 1789, an account of the appearances of Saturn's ring as seen from the planet, and a remark on the grey equatorial band on Saturn. The authors suggest that it may be a water-like liquid connected with the attraction exerted by the ring. But we cannot agree with the claim that, in the regions where the ring is at the zenith, it must \emph{necessarily} produce a constant tide many thousands of times greater than the tide that occurs on Earth. In fact, whether such a constant tide exists depends essentially on the shape of the solid part of the planet itself. If, as we must consider likely by analogy, that shape formed under the influence of all the forces belonging to Saturn's system, including those attributable to the attraction of the ring, then strictly speaking there can be no question of a constant tide at all.
VII. Jupiter. New determination of its rotation period, relative brightness of its satellites, and measurements of its equatorial and polar diameters.
VIII. Mars, especially its white polar spots.
IX. Venus. Eighteen drawings of the shape of the terminator really show only how difficult it is to ascertain anything precise about the rotation period, for which the values given by \textsc{Cassini} and \textsc{Bianchini} differ so enormously. The authors nevertheless see in those observations, however little they can settle the question, some indication in favour of the shorter rotation period. In the case of a peculiar radiant phenomenon that Mr. \textsc{Mädler} noticed on Venus on 7 April 1833, the \SourceLeaf{659}information reported affords no basis for a judgment.
X. Mercury. Measurement of its diameter during its transit in front of the Sun on 4 May 1832.
XI. Heliometric measurements of some double stars.
